\chapter{Introduction -- Giới thiệu}\label{ch:introduction}
\section{Bối cảnh và động lực nghiên cứu}
Dịch vụ dựa trên vị trí (Location-based Services, LBS) sử dụng thông tin không gian để hỗ trợ những
nhu cầu như tìm quán ăn, trạm sạc, cơ sở y tế hoặc địa điểm phù hợp trên đường đi. Khi người dùng di
chuyển, một yêu cầu đơn lẻ thường không đủ: thiết bị tiếp tục hỏi dữ liệu mới, nhận phản hồi và cập
nhật đáp án. Chuỗi tương tác này tạo ra một transcript có tọa độ, thứ tự và thời gian. Nhà cung cấp
không chỉ quan sát nơi một yêu cầu được gửi, mà còn có thể kết hợp nhiều yêu cầu với bản đồ và kiến
thức sẵn có về chuyển động.

Trajectory có tính tương quan: các điểm thuộc cùng một hành trình phải phù hợp với mạng đường,
thời gian di chuyển và hoạt động của người dùng. Một điểm chưa đủ rõ có thể trở nên dễ suy luận khi
đặt cạnh các điểm trước và sau nó. Nghiên cứu về tính duy nhất của dấu vết di chuyển cũng cho thấy
thông tin không gian--thời gian có thể mang tín hiệu nhận dạng mạnh trong một số tập dữ liệu
\cite{mobilityunique}. Kết quả đó tạo động lực cho nghiên cứu riêng tư, nhưng không được sử dụng để
mặc định rằng mọi trajectory tổng hợp đều đại diện đầy đủ cho danh tính người thật.

Trong bối cảnh này, thêm một điểm nhiễu là một thao tác kỹ thuật; điều cần giải thích là mức thông
tin mà \emph{cả quá trình sử dụng dịch vụ} có thể tiết lộ. Geo-indistinguishability (Geo-I) đưa ra một
cách phát biểu bảo đảm theo xác suất và khoảng cách vị trí \cite{geoi}. Tuy nhiên, một hệ thống online
còn có các quyết định tạo mới/tái sử dụng thông tin, chọn điểm truy vấn, gửi nhiều khung và sử dụng
cache. Muốn đưa ra bảo đảm cho hệ thống, cần phân tích những quyết định đó trong một mô hình chung.

\section{Các vấn đề nghiên cứu}
\subsection{Từ snapshot đến transcript nhiều bước}
Một cơ chế có thể tạo đầu ra khó phân biệt ở từng lần gửi nhưng vẫn tích lũy thông tin khi được gọi
nhiều lần. Attacker có thể khai thác chuỗi vị trí, tốc độ và quan hệ reachability. Quyết định cập nhật
hay giữ một điểm tham chiếu cũng dựa trên GPS; bỏ qua quyết định đó trong accounting có thể làm cận
của toàn chuỗi không còn đúng. Lập luận privacy vì vậy phải xác định lịch quan sát, lịch có thể đọc,
trạng thái được giữ qua phiên và cách hợp thành các bước.

Privacy được phát biểu bằng giới hạn mức transcript làm thay đổi khả năng phân biệt/suy luận.
Prior mạnh hoặc định danh trực tiếp vẫn có thể cho attacker nhiều thông tin từ trước.

\subsection{Chất lượng dịch vụ phụ thuộc đáp án POI}
Khoảng cách giữa điểm thật và điểm công bố chưa thể hiện đầy đủ trải nghiệm dịch vụ. Một vị trí giả
khá gần GPS vẫn có thể lấy thiếu địa điểm đúng khi có đường một chiều, sông hoặc quy tắc xếp hạng
khác nhau. Ngược lại, nhiều query locations và cache có thể tạo ra đủ ứng viên để thiết bị tìm được
đáp án tốt dù từng tọa độ công bố ở xa. Đánh giá utility cần đi qua chuỗi truy hồi POI, hợp/cache,
lọc và sắp xếp ở vị trí thật.

Nhu cầu tìm POI cũng có nhiều dạng: gần nhất, nhanh nhất, trong một bán kính và ít vòng đường trên
đường tới một đích đã biết tại thiết bị. Một dịch vụ có thể giữ nhu cầu này ở local và gửi một dạng
request chung. Việc không gửi payload thật phải được phân biệt với việc che mọi ý định tương quan
với tuyến đường hoặc thời điểm bắt đầu sử dụng dịch vụ.

\subsection{Khó khăn khi đối chiếu các nghiên cứu}
Các công trình liên quan có thể công bố một tọa độ thay thế, một tập chứa GPS thật và các dummy,
chèn fake packets, hoặc phát hành một bộ trajectory offline. Mỗi output contract làm thay đổi điều
attacker phải suy luận và cách tính metrics. Chẳng hạn, một metric chọn thành viên thật trong tập
không áp dụng nguyên dạng khi cơ chế không đưa GPS thật vào tập Q. JSD của dữ liệu tổng hợp cũng
không đo cùng tác vụ với lỗi định vị online \cite{semantic,fake,dpfetc,sok}.

Luận văn không khẳng định toàn lĩnh vực không có dataset hay metrics chung. Nhận định có giới hạn
là những công trình được khảo sát chưa cung cấp một protocol trực tiếp dùng chung cho các output và
mục tiêu đang xét. Vì vậy, cần xây dựng workload có kiểm soát, khai báo attacker và giữ hai lớp đo:
metrics chung theo tác vụ, cùng metrics gốc khi đầu vào/định nghĩa tương thích. Những metrics quen
thuộc như MAE hay Recall không được gọi là công thức mới của luận văn.

\section{Hệ thống hóa mục tiêu bảo vệ}
Luận văn sử dụng mười scenario để làm rõ \emph{target suy luận}, thay vì xem mọi nguy cơ là một score
privacy duy nhất. S1--S3 xét vị trí hiện tại, nơi dừng và đường đã đi; S4 xét liên kết người/thiết bị/
phương tiện qua các quan sát; S5/S6 xét cạnh kế tiếp và đích tương lai; S7 xét nội dung hoặc ý định
truy vấn; S8 xét thông tin tương quan từ người đi cùng; S9/S10 xét điểm đầu và cuối phiên.

S1--S3/S9--S10 là thực nghiệm chính; S4--S8 được đánh giá theo target và dữ liệu thực sự có.
Proxy nhóm tuyến và hai phiên căn clock được tách khỏi danh tính người/xe và quan sát đồng hành thực tế.

Workload SUMO và catalogue POI tạo điều kiện giữ cùng bản đồ, lịch và đáp án tham chiếu giữa các
phương pháp. Các giả thiết sinh dữ liệu và việc chia family cần được khai báo để hiểu các con số.
Chi tiết construction, metrics và coverage được đặt ở Chương~\ref{ch:experiments}, sau khi bài toán
và phương pháp đã được đặc tả.

\section{Hướng giải quyết và câu hỏi nghiên cứu}
GeoI-BQ v1 duy trì một điểm tham chiếu Z bằng cơ chế trên mạng đường dựa trên Geo-I và một phép kiểm
tra tái sử dụng có nhiễu. Thông tin đã bảo vệ được đưa vào một belief gần đúng để đánh giá các vị trí
có thể của người dùng. Belief điều chỉnh phân phối chọn \emph{đoạn Q} từ một thư viện chung. Đoạn được
chọn một lần rồi phát từng khung, giúp nối thiết kế online với accounting privacy rõ ràng. POI trả về
được hợp/cache và xử lý theo nhu cầu thật tại thiết bị.

Hai bước Z và Q là một lựa chọn thiết kế có phân vai: Z quản lý thông tin vị trí đã bảo vệ theo thời
gian; Q là hành động truy hồi dịch vụ mà server nhìn thấy. Tính an toàn không đến từ việc gọi đó là
``hai lần nhiễu'', mà từ các cận xác suất và hợp thành. Luận văn tập trung ba câu hỏi:
\begin{enumerate}[label=\textbf{RQ\arabic*.}]
\item Khi attacker quan sát toàn transcript trong một phiên, có thể chứng minh giới hạn nào cho
khả năng phân biệt trajectory và mức tăng suy luận?
\item Việc chọn Q, hợp/cache và local ranking bảo đảm utility ở mức nào; điều kiện nào cần để nối
điểm lập kế hoạch với chất lượng đáp án tại GPS thật?
\item Trong protocol được khai báo, GeoI-BQ đánh đổi privacy, utility và cost ra sao so với các
adaptation liên quan và các controls trên cùng workload?
\end{enumerate}

\section{Các đóng góp chính}
\begin{enumerate}
\item \textbf{Đặc tả bài toán và đánh giá theo contract.} Luận văn tổ chức S1--S10, phân biệt thông tin
riêng/công khai và quyền quan sát; xây dựng protocol cho những output khác nhau và giữ trường hợp
không áp dụng, failure, confidence interval cùng cost.
\item \textbf{Thiết kế phương pháp và phân tích toàn transcript.} Phương pháp nối REM/noisy reuse,
belief và kernel chọn chung đoạn Q có ngân sách tổng hữu hạn; phân tích privacy và utility nêu rõ
giả thiết cũng như quan hệ giữa bảo đảm lý tưởng và prototype.
\item \textbf{Thực nghiệm có thể tái lập.} So sánh chính dùng ba adaptation đã audit từ paper, kèm
Raw GPS và REM-Z; các phép đo mở rộng giữ target phù hợp và quá trình tăng cường attacker S8.
Bảng/hình được xuất từ artifact đã xác thực, không điều chỉnh defense để tìm kết quả thuận lợi.
\end{enumerate}

Geo-I, dự đoán riêng tư, hợp thành và hậu xử lý là kiến thức kế thừa \cite{geoi,predictive,dworkroth}.
Đóng góp của luận văn nằm ở đặc tả/cơ chế kết hợp và những phân tích cho policy đã nêu, không mặc định
là lần đầu có Geo-I, dummy hoặc local selection. Kết quả chính là một đánh đổi có điều kiện, không
là cam kết thắng mọi metrics hoặc bảo vệ mọi side channel.

\section{Tổ chức luận văn}
Chương~\ref{ch:related} phân tích nghiên cứu liên quan và khoảng trống đánh giá.
Chương~\ref{ch:method} đặc tả hệ thống, thuật toán và chứng minh privacy/utility trước thực nghiệm.
Chương~\ref{ch:experiments} trình bày metrics, dữ liệu, setup, kết quả chính và diagnostic bổ sung.
Chương~\ref{ch:summary} trả lời các câu hỏi nghiên cứu, tổng kết đóng góp, giới hạn và hướng tiếp theo.
Phụ lục giữ CIs, cấu hình/tái lập và các bảng giải thích nội bộ để phần chính tập trung vào lập luận.
